\documentclass[11pt]{article}
\usepackage[margin=1.1in]{geometry}
\usepackage{amsmath,amssymb,amsthm,mathtools,booktabs,array,tabularx,longtable}
\usepackage[dvipsnames]{xcolor}
\usepackage[colorlinks=true,linkcolor=MidnightBlue,citecolor=MidnightBlue,urlcolor=MidnightBlue]{hyperref}
\usepackage{fancyhdr}
\pagestyle{fancy}\fancyhf{}\renewcommand{\headrulewidth}{0pt}
\fancyfoot[R]{\thepage}
\fancyfoot[L]{\scriptsize\color{gray}Generated by the Transport Lab from approved parts of its store; not a deposited paper. Papers of record: \href{https://zenodo.org/communities/transport-papers}{Zenodo, Transport Papers}}
\fancypagestyle{plain}{\fancyhf{}\fancyfoot[R]{\thepage}\fancyfoot[L]{\scriptsize\color{gray}Generated by the Transport Lab from approved parts of its store; not a deposited paper. Papers of record: \href{https://zenodo.org/communities/transport-papers}{Zenodo, Transport Papers}}\renewcommand{\headrulewidth}{0pt}}
\newtheorem{theorem}{Theorem}[section]
\newtheorem{proposition}[theorem]{Proposition}
\newtheorem{lemma}[theorem]{Lemma}
\newtheorem{corollary}[theorem]{Corollary}
\newtheorem{conjecture}[theorem]{Conjecture}
\newtheorem{refutation}[theorem]{Refutation}
\newtheorem{prediction}[theorem]{Prediction}
\newtheorem{principle}[theorem]{Principle}
\newtheorem{observation}[theorem]{Observation}
\newtheorem{axiom}[theorem]{Axiom}
\theoremstyle{definition}
\newtheorem{definition}[theorem]{Definition}
\newtheorem{fact}[theorem]{Fact}
\newtheorem{assembled}[theorem]{Assembly}
\newtheorem{postulate}[theorem]{Postulate}
\newtheorem{premise}[theorem]{Premise}
\newtheorem{hypothesis}[theorem]{Hypothesis}
\newtheorem{remark}[theorem]{Remark}
\newtheorem{example}[theorem]{Example}
\theoremstyle{plain}
\setlength{\marginparwidth}{0.8in}
\newcommand{\changedmark}[1]{\marginpar{\raggedright\scriptsize\color{BrickRed}#1}}
\providecommand{\Aa}{}\renewcommand{\Aa}{\mathcal{A}}
\providecommand{\Ad}{}\renewcommand{\Ad}{\widetilde{\mathrm{Ad}}}
\providecommand{\Aev}{}\renewcommand{\Aev}{\mathcal{A}^{+}}
\providecommand{\Aod}{}\renewcommand{\Aod}{\mathcal{A}^{-}}
\providecommand{\B}{}\renewcommand{\B}{\mathrm{B}}
\providecommand{\C}{}\renewcommand{\C}{\mathbb{C}}
\providecommand{\CAR}{}\renewcommand{\CAR}{\mathrm{CAR}}
\providecommand{\CS}{}\renewcommand{\CS}{{\mathbb{S}^{+}}}
\providecommand{\CSm}{}\renewcommand{\CSm}{{\mathbb{S}^{-}}}
\providecommand{\Cl}{}\renewcommand{\Cl}{\mathrm{Cl}}
\providecommand{\Con}{}\renewcommand{\Con}{\mathcal{C}}
\providecommand{\Dl}{}\renewcommand{\Dl}{\Delta}
\providecommand{\EN}{}\renewcommand{\EN}{E_{N}}
\providecommand{\End}{}\renewcommand{\End}{\mathrm{End}}
\providecommand{\F}{}\renewcommand{\F}{\mathbb{F}}
\providecommand{\G}{}\renewcommand{\G}{\Gamma}
\providecommand{\HH}{}\renewcommand{\HH}{\mathbb{H}}
\providecommand{\IS}{}\renewcommand{\IS}{I\sigma_3}
\providecommand{\J}{}\renewcommand{\J}{\mathbf{J}}
\providecommand{\LH}{}\renewcommand{\LH}{L_{H}}
\providecommand{\Mat}{}\renewcommand{\Mat}{\mathrm{M}}
\providecommand{\Mo}{}\renewcommand{\Mo}{M^{\circ}}
\providecommand{\Mp}{}\renewcommand{\Mp}{M^{+}}
\providecommand{\Msun}{}\renewcommand{\Msun}{M_{\odot}}
\providecommand{\Mt}{}\renewcommand{\Mt}{\widetilde{M}}
\providecommand{\Ncone}{}\renewcommand{\Ncone}{\mathcal{N}}
\providecommand{\Neff}{}\renewcommand{\Neff}{N_{\mathrm{eff}}}
\providecommand{\PCS}{}\renewcommand{\PCS}{\mathbb{P}{\mathbb{S}^{+}}}
\providecommand{\PT}{}\renewcommand{\PT}{\mathbb{PT}}
\providecommand{\Pin}{}\renewcommand{\Pin}{\mathrm{Pin}}
\providecommand{\Q}{}\renewcommand{\Q}{\mathbb{Q}}
\providecommand{\R}{}\renewcommand{\R}{\mathbb{R}}
\providecommand{\Ric}{}\renewcommand{\Ric}{\mathrm{Ric}}
\providecommand{\SBH}{}\renewcommand{\SBH}{S_{\mathrm{BH}}}
\providecommand{\Sc}{}\renewcommand{\Sc}{\mathbb{S}}
\providecommand{\Sch}{}\renewcommand{\Sch}{\mathrm{P}}
\providecommand{\Sig}{}\renewcommand{\Sig}{\Sigma}
\providecommand{\Som}{}\renewcommand{\Som}{S_\omega}
\providecommand{\Sp}{}\renewcommand{\Sp}{\mathbb{S}}
\providecommand{\Sph}{}\renewcommand{\Sph}{\mathrm{S}}
\providecommand{\Spi}{}\renewcommand{\Spi}{S_\pi}
\providecommand{\Spin}{}\renewcommand{\Spin}{\mathrm{Spin}}
\providecommand{\Sth}{}\renewcommand{\Sth}{S_{\mathrm{th}}}
\providecommand{\T}{}\renewcommand{\T}{\mathbb{T}}
\providecommand{\TH}{}\renewcommand{\TH}{T_{\mathrm{H}}}
\providecommand{\Tr}{}\renewcommand{\Tr}{\operatorname{Tr}}
\providecommand{\UHF}{}\renewcommand{\UHF}{\mathrm{UHF}}
\providecommand{\Z}{}\renewcommand{\Z}{\mathbb{Z}}
\providecommand{\Zc}{}\renewcommand{\Zc}{\mathrm{Z}}
\providecommand{\cmg}{}\renewcommand{\cmg}{\mathrm{cm^{2}\,g^{-1}}}
\providecommand{\dS}{}\renewcommand{\dS}{\mathrm{dS}}
\providecommand{\dsc}{}\renewcommand{\dsc}{\mathrm{dsc}}
\providecommand{\dtPG}{}\renewcommand{\dtPG}{dt}
\providecommand{\dtS}{}\renewcommand{\dtS}{dt_{\mathrm{S}}}
\providecommand{\ellv}{}\renewcommand{\ellv}{\ell}
\providecommand{\er}{}\renewcommand{\er}{e_r}
\providecommand{\et}{}\renewcommand{\et}{e_t}
\providecommand{\ev}{}\renewcommand{\ev}{\mathrm{ev}}
\providecommand{\gl}{}\renewcommand{\gl}{\mathfrak{gl}}
\providecommand{\hatV}{}\renewcommand{\hatV}{\widehat{V}}
\providecommand{\hb}{}\renewcommand{\hb}{\bar{h}}
\providecommand{\hf}{}\renewcommand{\hf}{\hat{f}}
\providecommand{\id}{}\renewcommand{\id}{\mathrm{id}}
\providecommand{\kpc}{}\renewcommand{\kpc}{\mathrm{kpc}}
\providecommand{\lP}{}\renewcommand{\lP}{\ell_{P}}
\providecommand{\mD}{}\renewcommand{\mD}{\mathbf{m}_{D}}
\providecommand{\mb}{}\renewcommand{\mb}{\bar m}
\providecommand{\mnu}{}\renewcommand{\mnu}{\Sigma m_{\nu}}
\providecommand{\nInf}{}\renewcommand{\nInf}{n_\infty}
\providecommand{\nO}{}\renewcommand{\nO}{n_0}
\providecommand{\ninf}{}\renewcommand{\ninf}{n_{\infty}}
\providecommand{\nz}{}\renewcommand{\nz}{n_{0}}
\providecommand{\proj}{}\renewcommand{\proj}{\operatorname{proj}}
\providecommand{\pv}{}\renewcommand{\pv}{\pi}
\providecommand{\rad}{}\renewcommand{\rad}{\mathrm{rad}}
\providecommand{\sgn}{}\renewcommand{\sgn}{\operatorname{sgn}}
\providecommand{\shA}{}\renewcommand{\shA}{\mathcal{A}}
\providecommand{\shC}{}\renewcommand{\shC}{\mathcal{C}}
\providecommand{\shE}{}\renewcommand{\shE}{\mathcal{E}}
\providecommand{\shJ}{}\renewcommand{\shJ}{\mathcal{J}}
\providecommand{\shK}{}\renewcommand{\shK}{\mathcal{K}}
\providecommand{\shL}{}\renewcommand{\shL}{L}
\providecommand{\shW}{}\renewcommand{\shW}{\mathcal{W}}
\providecommand{\sll}{}\renewcommand{\sll}{\mathfrak{sl}}
\providecommand{\so}{}\renewcommand{\so}{\mathfrak{so}}
\providecommand{\su}{}\renewcommand{\su}{\mathfrak{su}}
\providecommand{\supp}{}\renewcommand{\supp}{\operatorname{supp}}
\providecommand{\tr}{}\renewcommand{\tr}{\operatorname{tr}}
\providecommand{\uu}{}\renewcommand{\uu}{\mathfrak{u}}
\providecommand{\br}{}\renewcommand{\br}[1]{\langle #1\rangle}
\providecommand{\bra}{}\renewcommand{\bra}[1]{\langle #1|}
\providecommand{\file}{}\renewcommand{\file}[1]{\texttt{#1}}
\providecommand{\hp}{}\renewcommand{\hp}[2]{( #1 \,|\, #2 )}
\providecommand{\ip}{}\renewcommand{\ip}[2]{\langle #1, #2\rangle}
\providecommand{\ket}{}\renewcommand{\ket}[1]{|#1\rangle}
\providecommand{\rel}{}\renewcommand{\rel}[2]{D\bigl(#1 \,\|\, #2\bigr)}
\providecommand{\rev}{}\renewcommand{\rev}[1]{\widetilde{#1}}
\providecommand{\slashed}{}\renewcommand{\slashed}[1]{{#1}\mkern-10mu/}
\title{\bfseries Einstein's equations from horizon thermodynamics}
\author{Leon Butler\\ \normalsize\href{mailto:lbutler.math@gmail.com}{\texttt{lbutler.math@gmail.com}}\\ \normalsize ORCID: \href{https://orcid.org/0009-0006-4449-5617}{0009-0006-4449-5617}}
\date{Generated by the Transport Lab, 27 October 2026}
\begin{document}
\maketitle
\begin{abstract}
\normalsize
This trail follows the route from the thermodynamics of horizons to Einstein's equations. Black-hole temperature and entropy alone cannot single out Einstein's theory, because a whole family of theories shares the same black holes.

Jacobson's argument closes the gap by applying the heat law $\delta Q = T\,\delta S$ to every small horizon that matter crosses. Under the series' local-horizon premise (C$_{\mathrm{loc}}$), the store proves Einstein's equations along that route, with $\Lambda$ an integration constant and $G_{\mathrm{cl}} = 1/(4\eta)$.

The remaining steps take the premise apart: its temperature clauses are one joint thermal condition, and its third clause adds nothing beyond them.
\end{abstract}
\newpage
\noindent \textbf{Authorship and disclosure. }The founding idea of the series is the author's: that division by zero in a Clifford algebra should be read not as a prohibition but as a transport into a higher-dimensional algebra, and that the transported division need not be accompanied by an inverse operation. The direction of investigation and the critical review of the results are likewise the author's. The mathematical development and the drafting of the text were carried out in extended collaboration with Claude Fable 5 (Anthropic), an AI system.

\noindent\fbox{\parbox{\dimexpr\linewidth-2\fboxsep-2\fboxrule}{\textbf{This is not a deposited paper.} It was generated by the Transport Lab on 27 October 2026 and assembled from approved parts of its store. It follows results across several papers of the series; each is stated with its proof and credited to the paper that first proved it. The official versions of those papers are on Zenodo, in the Transport Papers community, \url{https://zenodo.org/communities/transport-papers}.}}
\newpage

\section*{Introduction}
The walk first says why black holes alone are not enough, then sets out the two imported tools, then proves the equations, and ends by examining what the premise it rests on actually contains.

\section{What we assume}

\begin{postulate}[(C$_{\mathrm{loc}}$): on every local Rindler horizon the entropy of a cut is the integral of a curvature-independent area density (Loc) and, before matter crosses, the modes on the light rays are jointly thermal for the local boost (R$_{\mathrm{loc}}$)]\label{c:synthesis.c-loc-local-horizon-premise}
The series computes a temperature and an entropy for the edges of black holes and of the universe's horizon, using quantum modes it places on those edges. Jacobson showed in 1995 that Einstein's equations of gravity follow if every small piece of horizon around every point, including the horizon seen by an accelerating observer, has a temperature set by the acceleration and an entropy proportional to its area, by the same amount everywhere. This assumption says exactly that, with the quantum modes running continuously along the light rays of each small horizon and being the same field as the matter that falls through it. It has two parts: the entropy of a slice of the horizon is its area times a fixed number that does not depend on how strong the gravitational field is there, and before matter falls through, the modes are thermal at the accelerating observer's temperature. Where there is no gravitational field the second part now follows from a simpler assumption with no temperature in it; the first part is still assumed. With it, the store proves Einstein's equations, with the strength of gravity set by the entropy per area and the cosmological constant left free. The assumption would fail if the entropy per area turned out to depend on the strength of the gravitational field or on matter, or if gravity were found to deviate from Einstein's equations.
\end{postulate}

\section{Einstein's equations in the completion}

Each step either supplies an ingredient of the thermodynamic route to Einstein's equations or tests the premise that route needs.

We start from what the series already has: black holes with a temperature, a count of modes and an entropy.

Every Lagrangian \[ \frac{R}{16\pi G} + f($\mathrm{Ric}$) \] with $f$ at least quadratic in the Ricci tensor shares the Schwarzschild field, its temperature and, assuming Wald entropy, its entropy $A/4G$. So these inputs cannot pick out Einstein's equations; what would separate the theories is entropy on local horizons crossed by matter.

\begin{theorem}[The series' black-hole thermodynamics holds unchanged in every Ricci-polynomial extension of Einstein's theory, so on its own it cannot yield Einstein's equations]\label{c:synthesis.black-hole-data-do-not-fix-field-equations}
The series has three facts about black holes: their temperature, how many quantum modes they hold, and Hawking's argument that turns the temperature into an entropy. A whole family of theories of gravity, differing from Einstein's in how matter sources the gravitational field, has exactly the same black holes, with the same temperature and, if horizon entropy is given by the standard (Wald) rule, the same entropy. So those three facts, on their own, cannot single out Einstein's equations. The known derivation of Einstein's equations from thermodynamics works because it applies the heat law to tiny horizons that matter is flowing through. The series does not yet say what entropy such a horizon has, and that is the step it would need to supply.
\end{theorem}
\noindent\emph{Rests on:} \ref{c:synthesis.quarter-from-temperature}, \ref{c:scrambling.mode-count}, \ref{c:singularity.rate-discharge-hypothesis-iii}, \ref{c:synthesis.imported-wald-noether-charge-entropy}.

The first tool is geometry.

A bundle of light rays focuses as it travels: its expansion $\theta$ obeys \[ (k \cdot D)\theta = -\theta\^{}{2}/2 - \sigma\^{}{2} + \omega\^{}{2} - \Ric\^{}{\circ}(k,k) \] with $\sigma$ the shear and $\omega$ the twist. Matter and field strength along the rays pull them together, so the area of the bundle shrinks.

\begin{fact}[Imported: the focusing and shear equations for a null geodesic congruence (Raychaudhuri 1955, Sachs 1961)]\label{c:synthesis.imported-raychaudhuri-null-congruence}
Light rays travelling side by side can converge or spread out. Raychaudhuri and Sachs found the equations that say how a bundle of rays changes as it travels: matter and the gravitational field along the rays focus them together, and a separate part of the field strength squeezes the bundle out of shape. These equations are standard textbook geometry; the series uses them in the step where the energy flowing through a small horizon is converted into the strength of the gravitational field.
\emph{Source:} A. Raychaudhuri, Relativistic cosmology. I, Phys. Rev. 98 (1955), 1123--1126.; R. K. Sachs, Gravitational waves in general relativity. VI. The outgoing radiation condition, Proc. R. Soc. Lond. A 264 (1961), 309--338.; R. M. Wald, General Relativity, University of Chicago Press (1984)..
\end{fact}

The second tool is thermodynamics.

If every local Rindler horizon has the Unruh temperature, an entropy $\eta$ per unit area, and obeys $\delta Q = T\,\delta S$ for the matter's boost-energy flux, then the gravitational field obeys \[ $\mathcal{G}$(a) + \Lambda a = -(2\pi/\eta)\,T(a) \] with $\Lambda$ undetermined. This is Jacobson's 1995 result, written in the store's sign.

\begin{fact}[Imported: Einstein's equation as an equation of state (Jacobson 1995)]\label{c:synthesis.imported-jacobson-clausius-1995}
In 1995 Ted Jacobson showed that Einstein's equations can be read as a law of thermodynamics. If every small patch of horizon, as seen by an accelerating observer, has a temperature set by the acceleration and an entropy proportional to its area, and if heat crossing it changes its entropy by heat divided by temperature, then the gravitational field must obey exactly Einstein's equations, read here as the equations of a field on flat spacetime. The strength of gravity is fixed by the entropy per area, and the cosmological constant is left free. This is a known result taken from the literature.
\emph{Source:} T. Jacobson, Thermodynamics of spacetime: the Einstein equation of state, Phys. Rev. Lett. 75 (1995), 1260--1263..
\end{fact}

The series' local-horizon premise (C$_{\mathrm{loc}}$) supplies the temperature and the entropy, and this step carries the argument through inside the store.

Every further hypothesis is stated: Clausius with $\delta S = \eta\,\delta A$, a thin patch with vanishing expansion and shear, a conserved $T_{ab}$, a $C^{3}$ derived metric, four dimensions. Then \[ $\mathcal{R}$(a) + 8\pi G\_{$\mathrm{cl}$}\,T(a) = (\Lambda + 4\pi G\_{$\mathrm{cl}$}\,\tau)\,a \] with $\Lambda$ an integration constant and $G_{\mathrm{cl}} = 1/(4\eta)$.

\begin{theorem}[Under (C$_{\mathrm{loc}}$) and the Clausius relation, Jacobson's argument gives Einstein's equations with $\Lambda$ an integration constant and $G_{\mathrm{cl}} = 1/(4\eta)$; every further hypothesis named]\label{c:synthesis.einstein-equations-from-c-loc-and-clausius}
In 1995 Ted Jacobson showed that Einstein's equations of gravity follow if every small patch of horizon, as seen by an accelerating observer, has a temperature set by the acceleration and an entropy proportional to its area, and if heat crossing the patch raises its entropy by heat divided by temperature. The series assumes, as a premise, that its own horizon quantum modes give that temperature and that entropy on all such small horizons. This result writes the argument out step by step under that premise and checks each algebraic step exactly. It also lists everything else the argument needs: that the heat is the energy carried by matter and nothing else changes the entropy, that the horizon patch is momentarily neither expanding nor shearing, that matter is described by an ordinary conserved energy tensor, that the gravitational field is smooth enough, and that space-time has four dimensions. Given all that, Einstein's equations follow, with the strength of gravity fixed by the entropy per area and the cosmological constant left free. What singles out Einstein's equations, rather than more complicated versions of them, is that the entropy per area is the same everywhere and does not depend on the strength of the gravitational field there. The derivation is conditional: it shows what the premise would deliver, not that the premise holds.
\end{theorem}
\noindent\emph{Rests on:} \ref{c:synthesis.imported-jacobson-clausius-1995}, \ref{c:synthesis.einstein-condition-is-infinity-vacuum-incidence}, \ref{c:synthesis.imported-raychaudhuri-null-congruence}, \ref{c:synthesis.c-loc-local-horizon-premise}.

The derivation is conditional, so the rest of the walk asks what the premise contains.

Its temperature clauses are one condition: joint KMS for the local boost at rapidity period $2\pi$, a new premise stronger than the local form of (R). Derived is only the geometric half, that this period is where the local boost's Euclidean section is regular, giving \[ T = a/2\pi \] up to a relative correction of order $\rho^{2}$.

\begin{theorem}[On local Rindler horizons and small diamonds, clauses (ii) and (iii) of (C$_{\mathrm{loc}}$), with (iii) read on the patch's whole mode algebra and the boost acting on the modes by phases, are one condition, joint KMS for the local boost at rapidity period $2\pi$, the period at which the local boost's Euclidean section is regular at leading order]\label{c:synthesis.local-boost-temperature-is-local-regularity}
The assumption that every small horizon is hot at the Unruh temperature and that its quantum state's internal clock runs along the horizon's own boost was stated as two separate requirements. This result shows that, read as statements about all the modes of the horizon together, they are one. The geometry part can be made local: near any point, however strong the gravitational field is, the smoothness condition that sets the temperature at a black hole gives the Unruh temperature, up to small corrections. For a quantum state on a finite set of modes, being thermal for a flow and having that flow as its internal clock are the same property. What is left is one assumption, that the modes of every small horizon are jointly thermal at that temperature. At black holes the series makes the corresponding assumption in two parts, one for each mode separately and one for how the modes combine; the local version needs both, so it is a new assumption, stronger than the black-hole smoothness rule alone. It is not derived here from anything weaker.
\end{theorem}
\noindent\emph{Rests on:} \ref{c:singularity.euclidean-section-period}, \ref{c:angle.imported-theorems-deliver-their-hypotheses}, \ref{c:singularity.kms-uniqueness-factor}, \ref{c:hawking.gibbs-states-are-kms-conversely-finite}, \ref{c:singularity.modular-theory-matrix-factor-proved-outright}, \ref{c:hawking.kms-restricts-invariant-subalgebras}, \ref{c:hawking.imported-kms-passes-inductive-limits}, \ref{c:hawking.full-uniqueness-given-dynamics}, \ref{c:singularity.rate-rigidity-temperature-not-algebra-intrinsic}.

Next the entropy clause.

If the entropy density of a small patch depends only on its flat tangent data, not on the field strength or matter at the point, the Lorentz group forces one universal constant $\eta$. Calibrated on the Schwarzschild family under (P), (R) and Hawking's hypothesis, \[ \eta = \tfrac14 \] with no mode count. That independence from curvature excludes the Wald-type entropies of the rival theories met at the start.

\begin{theorem}[If the entropy density of a small horizon patch depends only on its flat tangent data (null ray and edge plane), and not on the curvature, matter or boost normalisation at the point, it is one universal constant; calibrated on the Schwarzschild family under (P), (R) and Hawking's hypothesis, with the hole's entropy taken to be the integral of that density, it is 1/4 with no count]\label{c:synthesis.c-loc-entropy-clause-curvature-independent-density}
The local-horizon assumption gives every small horizon an entropy proportional to its area, with the same constant everywhere. This result shows that "the same constant everywhere" is automatic once one assumes that the entropy per area of a small piece of horizon depends only on how the piece sits among the flat directions at its point, and not on how strong the gravitational field is there or what matter is present. That second part is a real assumption: theories of gravity with extra curvature terms give entropies that do depend on the strength of the gravitational field, and this assumption is what rules them out. Given it, a change of viewpoint or of the observer's acceleration turns any small horizon into any other without changing its area, so the entropy per area can only be one number. That number is then fixed by black holes: if a black hole's entropy is the sum of this density over its horizon, and under the series' assumptions about its temperature and Hawking's thermodynamic reasoning, the constant is one quarter, with no need to count modes. The series' mode count gives the same quarter only at the value of its calibration constant that was chosen to match black-hole entropy, as already known.
\end{theorem}
\noindent\emph{Rests on:} \ref{c:synthesis.imported-jacobson-clausius-1995}, \ref{c:synthesis.quarter-from-temperature}, \ref{c:synthesis.de-sitter-entropy-pi-n-squared}, \ref{c:synthesis.imported-kerr-surface-gravity-entropy}, \ref{c:synthesis.clausius-route-newton-constant-is-residue}.

Putting modes on every horizon raises a bookkeeping question, because the cutoff premises place a black hole's modes at a locus inside its horizon.

If the two placements are independent and the hole's entropy is the total over all its modes, Jacobson's constant becomes $G_{\mathrm{cl}} = 2G$ at the area-law calibration. If they are identified, they inherit the literal Kerr count's conflict with $A_{+}/4$, which counts against the cutoff placement on spinning holes, not against (C$_{\mathrm{loc}}$).

\begin{theorem}[(C$_{\mathrm{loc}}$)(i) is a second placement of the count, not an extension of (C1) and (C2); at a hole's horizon, independent placements double Newton's constant under the total-entropy reading, and the literal Kerr count's conflict with $A_{+}/4$, present in every reading, counts against the cutoff placement, not against (C$_{\mathrm{loc}}$)]\label{c:synthesis.c-loc-is-a-second-placement-of-the-count}
The series places a black hole's hidden quantum modes deep inside the hole, where the curvature reaches the Planck scale, and their number matches the hole's known entropy. A newer assumption, needed to derive Einstein's equations, puts modes with the same entropy per unit area on every small horizon, including the horizon of a black hole itself. This result shows that the two placements never touch: the deep locus lies inside the horizon of every hole heavier than half a Planck mass. So a black hole would have modes in two places. If the two sets are independent and the hole's entropy is the total over both, the entropy is counted twice, and fixing that would make the derived strength of gravity twice the true value. If they are the same set, then for fast-spinning holes the deep count disagrees with the known entropy, as it already does without the new assumption. That disagreement appears whichever way the two sets are related, so it counts against placing the modes deep inside spinning holes, not against the new assumption. The assumption still needs a rule saying how the two sets are related.
\end{theorem}
\noindent\emph{Rests on:} \ref{c:singularity.stretched-singularity-count}, \ref{c:scrambling.epoch-modular-period-two-unknowns-become}, \ref{c:scrambling.area-law-test}, \ref{c:synthesis.clausius-route-newton-constant-is-residue}, \ref{c:synthesis.imported-jacobson-clausius-1995}, \ref{c:synthesis.kerr-ring-count-misses-area-law}, \ref{c:synthesis.imported-kerr-surface-gravity-entropy}, \ref{c:synthesis.quarter-from-temperature}.

The last step looks at the third clause as first written, that the modes' modular flow acts geometrically on nested sub-patches.

A boost acting by phases is almost periodic, so it maps no closed subspace properly into itself and cannot push a smaller patch into a larger one. The clause is consistent only as the boost dynamics on the whole patch, where for a faithful state it restates the temperature clause.

\begin{theorem}[A boost acting by phases moves no C$^{*}$-subalgebra, and in a GNS representation no weakly closed subalgebra, properly into itself, so clause (iii) of (C$_{\mathrm{loc}}$) cannot act geometrically on nested sub-patches of the mode algebra its clause (i) cites]\label{c:synthesis.c-loc-modular-clause-moves-no-nested-patch}
The new assumption needed for Einstein's equations has three parts: small horizons carry the series' quantum modes, those modes are at the right temperature, and their natural time flow is the geometric flow of the horizon. This result shows that the third part can only be meant in a weak sense. The series' modes evolve by simple rotations of phases, and such an evolution always comes back close to where it started. A genuinely geometric flow would instead push the modes of a smaller piece of horizon ever deeper into a larger piece, never returning. The two cannot both hold. So the third part must be read as a statement about the modes of the whole patch, where it adds nothing beyond the temperature part; the stronger geometric reading would need a continuous family of modes, which the assumption does not name.
\end{theorem}
\noindent\emph{Rests on:} \ref{c:synthesis.point-spectrum-boost-admits-no-half-sided-translations}.

\section{Results used from elsewhere}
The proofs above use the following results of the series, stated briefly; each is proved in the paper named.

\begin{fact}[Imported: the theorems that deliver $2\pi$, and their hypotheses]\label{c:angle.imported-theorems-deliver-their-hypotheses}
Under the spectrum condition, covariance and locality, the modular flow of a wedge or half-sided algebra is the boost at rate $2\pi$, and a regular stationary quasi-free state at a bifurcate Killing horizon is KMS at $\kappa/2\pi$.
\emph{Source:} J. J. Bisognano and E. H. Wichmann, On the duality condition for a Hermitian scalar field, J. Math. Phys. 16 (1975), 985--1007.; J. J. Bisognano and E. H. Wichmann, On the duality condition for quantum fields, J. Math. Phys. 17 (1976), 303--321.; H.-J. Borchers, The CPT-theorem in two-dimensional theories of local observables, Commun. Math. Phys. 143 (1992), 315--332.; P. D. Hislop and R. Longo, Modular structure of the local algebras associated with the free massless scalar field theory, Commun. Math. Phys. 84 (1982), 71--85.; R. Brunetti, D. Guido and R. Longo, Modular structure and duality in conformal quantum field theory, Commun. Math. Phys. 156 (1993), 201--219.; G. L. Sewell, Quantum fields on manifolds: PCT and gravitationally induced thermal states, Ann. Phys. 141 (1982), 201--224.; B. S. Kay and R. M. Wald, Theorems on the uniqueness and thermal properties of stationary, nonsingular, quasifree states on spacetimes with a bifurcate Killing horizon, Phys. Rep. 207 (1991), 49--136.; W. G. Unruh, Notes on black-hole evaporation, Phys. Rev. D 14 (1976), 870--892..
\end{fact}

\begin{fact}[Full uniqueness, given the dynamics]\label{c:hawking.full-uniqueness-given-dynamics}
Given the boost dynamics and $\beta$, the I-state is the unique $\tau$-KMS state of the boundary-mode algebra among all states.
\emph{Source:} T7, Cor 3.8 (\href{https://doi.org/10.5281/zenodo.22722121}{doi:10.5281/zenodo.22722121}).
\end{fact}

\begin{fact}[Gibbs states are KMS, and conversely, in finite dimensions]\label{c:hawking.gibbs-states-are-kms-conversely-finite}
On a full matrix algebra with dynamics generated by a diagonal Hamiltonian, the Gibbs state at inverse temperature $\beta$ is KMS, and it is the unique KMS state at that $\beta$, degenerate spectrum included.
\emph{Source:} T7, Lem 3.1 (\href{https://doi.org/10.5281/zenodo.22722121}{doi:10.5281/zenodo.22722121}).
\end{fact}

\begin{fact}[Imported: KMS passes to inductive limits]\label{c:hawking.imported-kms-passes-inductive-limits}
A state that is KMS at $\beta$ on each subalgebra of an increasing chain of dynamics-invariant subalgebras is KMS at $\beta$ on the closure of their union.
\emph{Source:} O. Bratteli and D. W. Robinson, Operator Algebras and Quantum Statistical Mechanics 2, 2nd ed., Springer, Berlin (1997)., Prop.~5.3.25.
\end{fact}

\begin{fact}[KMS restricts to invariant subalgebras]\label{c:hawking.kms-restricts-invariant-subalgebras}
A KMS state restricts to a KMS state at the same inverse temperature on any unital subalgebra invariant under the dynamics.
\emph{Source:} T7, Lem 3.2 (\href{https://doi.org/10.5281/zenodo.22722121}{doi:10.5281/zenodo.22722121}).
\end{fact}

\begin{fact}[The area-law test]\label{c:scrambling.area-law-test}
The mode count equals $\SBH = 4\pi M^2$ if and only if $\nu\,\Delta t = M/2$ (in the 1-nat convention; $M/(2\ln 2)$ when each mode is credited its $\ln 2$), an epoch of order one Schwarzschild time.
\emph{Source:} T8, Prop 3.5 (\href{https://doi.org/10.5281/zenodo.22722176}{doi:10.5281/zenodo.22722176}).
\end{fact}

\begin{fact}[The epoch is the modular period; two unknowns become one]\label{c:scrambling.epoch-modular-period-two-unknowns-become}
Under the intrinsic-clock premise the counted epoch is the modular period $\beta = 8\pi M$, turning the area law into the prediction $\nu = 1/16\pi$ in the 1-nat convention.
\emph{Source:} T8, Thm 4.3 (\href{https://doi.org/10.5281/zenodo.22722176}{doi:10.5281/zenodo.22722176}).
\end{fact}

\begin{fact}[The mode count]\label{c:scrambling.mode-count}
To leading order in $r_c/2M$, the boundary-mode count over a Killing-time span $\Delta t$ is $N = 8\pi\nu M\,\Delta t/\lP^2$, accruing at the rate $\nu/\TH$.
\emph{Source:} T8, Thm 3.4 (\href{https://doi.org/10.5281/zenodo.22722176}{doi:10.5281/zenodo.22722176}).
\end{fact}

\begin{fact}[The Euclidean section, and its period]\label{c:singularity.euclidean-section-period}
The Euclidean section of the Schwarzschild field is smooth at its tip, the Euclidean counterpart of the bifurcation sphere, if and only if imaginary Killing time has period $p = 2\pi/\kappa = 8\pi M$; any other period leaves a conical point of total angle $\kappa p$.
\emph{Source:} T5, Prop 7.18 (\href{https://doi.org/10.5281/zenodo.22721790}{doi:10.5281/zenodo.22721790}).
\end{fact}

\begin{fact}[KMS uniqueness at a factor]\label{c:singularity.kms-uniqueness-factor}
On a matrix factor, for each nonzero inverse temperature the only KMS state of an inner flow is the Gibbs state of its generator.
\emph{Source:} T5, Thm 7.13 (\href{https://doi.org/10.5281/zenodo.22721790}{doi:10.5281/zenodo.22721790}).
\end{fact}

\begin{fact}[Modular theory of a matrix factor, proved outright]\label{c:singularity.modular-theory-matrix-factor-proved-outright}
For a full matrix algebra with a cyclic separating vector, the modular operator, modular conjugation and modular flow are computed explicitly, and the state is KMS for that flow.
\emph{Source:} T5, Lem 7.1 (\href{https://doi.org/10.5281/zenodo.22721790}{doi:10.5281/zenodo.22721790}).
\end{fact}

\begin{fact}[The rate, and the discharge of hypothesis (iii)]\label{c:singularity.rate-discharge-hypothesis-iii}
Under (P) and (R) every mode's marginal is the Gibbs state at the Hawking temperature $T = 1/8\pi M$, so each marginal is faithful and the cyclic-and-separating hypothesis of modular descent is derived rather than assumed.
\emph{Source:} T5, Thm 7.20 (\href{https://doi.org/10.5281/zenodo.22721790}{doi:10.5281/zenodo.22721790}).
\end{fact}

\begin{fact}[Rate rigidity: the temperature is not algebra intrinsic]\label{c:singularity.rate-rigidity-temperature-not-algebra-intrinsic}
Automorphisms of the single-mode algebra that intertwine the boost's phase flow $\operatorname{Ad}(e^{itK})$ move its KMS states only by $\beta \mapsto \pm\beta$, so the value of the temperature depends on a parametrisation of the flow that the algebra does not supply.
\emph{Source:} T5, Thm 7.16 (\href{https://doi.org/10.5281/zenodo.22721790}{doi:10.5281/zenodo.22721790}).
\end{fact}

\begin{fact}[The stretched singularity, and the count]\label{c:singularity.stretched-singularity-count}
Under (C1) and (C2), the stretched singularity's volume per unit Killing time is $8\pi M$ at leading order, so it gains about $\nu/\TH$ modes per unit time.
\emph{Source:} T5, Prop 10.6 (\href{https://doi.org/10.5281/zenodo.22721790}{doi:10.5281/zenodo.22721790}).
\end{fact}

\begin{fact}[On Jacobson's Clausius route the epoch count returns Newton's constant only through the area-law condition, a consistency; the missing step is a count on non-Killing horizons]\label{c:synthesis.clausius-route-newton-constant-is-residue}
On Jacobson's Clausius route (the imported theorem (J)) the epoch count ($4\pi\nu$ modes per Planck area per period under (IC); the accrual count fails Clausius) would give, if it held on every local Rindler horizon, $G_{\mathrm{cl}} = G/(16\pi\nu\langle s\rangle)$, and $G_{\mathrm{cl}} = G$ is the area-law condition that already defines the registered $\nu$: a consistency, not a determination of $G$; the missing step is a count on non-Killing horizons.
\emph{Source:} Transport Lab claim synthesis.clausius-route-newton-constant-is-residue, version 4.
\end{fact}

\begin{fact}[The cosmological horizon's entropy is pi N\^{}2]\label{c:synthesis.de-sitter-entropy-pi-n-squared}
Under the imported Einstein-Hilbert area law, the cosmological horizon's entropy is exactly $\pi N^{2} = \pi M_{P}^{2}/Q(I)$ (and $\pi(N_{w}^{2} + 9/4)$ with $N_{w} = \operatorname{Im} w(M_{P})$).
\emph{Source:} Transport Lab claim synthesis.de-sitter-entropy-pi-n-squared, version 5.
\end{fact}

\begin{fact}[The standard null-cone form of the Einstein condition, translated into vacuum language: infinity vacua unexcited along their own light rays]\label{c:synthesis.einstein-condition-is-infinity-vacuum-incidence}
A translation, not a derivation: the standard equivalence ``Einstein if and only if $\Sch(\ell,\ell) = 0$ for all null $\ell$'' reads, through the Fock vacua annihilated by the pure-scale tractor, as ``every infinity vacuum is unexcited along its own light ray'' ($\Ric = (R/4)g$, $R$ then constant by the contracted Bianchi identity), with ``unexcited in every direction'' giving $\Ric = 0$, one unexcited vacuum giving $\Ric = 2c\,\ell\otimes\ell$, and, for a conserved source $T$, own-ray classes $4\pi G\,T(\ell,\ell)$ giving the Einstein equations with $\Lambda$ an integration constant; point vacua are unexcited along their own rays in every metric.
\emph{Source:} Transport Lab claim synthesis.einstein-condition-is-infinity-vacuum-incidence, version 3.
\end{fact}

\begin{fact}[Imported: the surface gravity and the Bekenstein-Hawking entropy of the Kerr horizon]\label{c:synthesis.imported-kerr-surface-gravity-entropy}
The Kerr field's horizon locus has surface gravity $\kappa_{+} = (r_{+} - r_{-})/(2(r_{+}^{2} + a^{2}))$, cross-section area $A_{+} = 4\pi(r_{+}^{2} + a^{2})$ and Bekenstein-Hawking entropy $A_{+}/4$.
\emph{Source:} J. M. Bardeen, B. Carter and S. W. Hawking, The four laws of black hole mechanics, Commun. Math. Phys. 31 (1973), 161--170.; S. W. Hawking, Particle creation by black holes, Commun. Math. Phys. 43 (1975), 199--220.; J. D. Bekenstein, Black holes and entropy, Phys. Rev. D 7 (1973), 2333--2346.; B. Carter, Global structure of the Kerr family of gravitational fields, Phys. Rev. 174 (1968), 1559--1571..
\end{fact}

\begin{fact}[Imported: Wald's Noether-charge entropy and its Jacobson-Myers form for a Gauss-Bonnet term]\label{c:synthesis.imported-wald-noether-charge-entropy}
Wald's entropy, $2\pi$ times the horizon Noether charge, $-2\pi\oint E^{abcd}\epsilon_{ab}\epsilon_{cd}\,dA$ on a cross-section of the horizon locus, obeys the first law of stationary black holes, and for $(R - 2\Lambda + \alpha E_{4})/16\pi G$ it is the Jacobson-Myers functional $\frac{1}{4G}\oint(1 + 2\alpha\,{}^{(2)}\!R)\,dA$.
\emph{Source:} R. M. Wald, Black hole entropy is Noether charge, Phys. Rev. D 48 (1993), R3427--R3431.; V. Iyer and R. M. Wald, Some properties of Noether charge and a proposal for dynamical black hole entropy, Phys. Rev. D 50 (1994), 846--864.; T. Jacobson and R. C. Myers, Black hole entropy and higher curvature interactions, Phys. Rev. Lett. 70 (1993), 3684--3687., Eq. (4); T. Jacobson, G. Kang and R. C. Myers, On black hole entropy, Phys. Rev. D 49 (1994), 6587--6598..
\end{fact}

\begin{fact}[For Kerr with a >= 2(2M)\^{}{2/3} the curvature-cutoff locus is a timelike tube beyond the inner horizon, and the counting rule read literally misses the area law by a factor of order M\^{}{-1/3}]\label{c:synthesis.kerr-ring-count-misses-area-law}
For Kerr, if $a \ge 2(2M)^{2/3}$, the Planck-curvature locus $|z| = (2M)^{1/3}$ is a timelike tube beyond the inner horizon with volume rate $4\pi^2(2M)^{2/3}(1 - \eta)$, so the counting rule read literally gives $N/(A_+/4) = (\pi/2)(2M)^{2/3}(1 - \eta)/\sqrt{M^2 - a^2}$, of order $M^{-1/3}$, against $1$ for Schwarzschild.
\emph{Source:} Transport Lab claim synthesis.kerr-ring-count-misses-area-law, version 5.
\end{fact}

\begin{fact}[On a mode space where the boost acts by phases, it dilates no translation group, so the Borchers route to (R) needs the continuum]\label{c:synthesis.point-spectrum-boost-admits-no-half-sided-translations}
A boost that acts by phases on a generating set of modes (pure point spectrum: windows, the countable mode space, the array of (Q)) dilates no strongly continuous translation group, even on the GNS space of a KMS state, so the Borchers and Bisognano-Wichmann route to (R) must be posed on the continuum mode space.
\emph{Source:} Transport Lab claim synthesis.point-spectrum-boost-admits-no-half-sided-translations, version 3.
\end{fact}

\begin{fact}[The quarter from the series' temperature, by Hawking's argument]\label{c:synthesis.quarter-from-temperature}
Under (P), (R) and the Clausius relation for the hole (Hawking's argument), $S = \int dM/T = 4\pi M^{2} = A/4$: the quarter is fixed by the period of (R).
\emph{Source:} Transport Lab claim synthesis.quarter-from-temperature, version 1.
\end{fact}

\section*{Conclusion}
The thread runs from black holes that cannot choose a theory of gravity, through the focusing of light and the heat law on small horizons, to Einstein's equations proved under one premise.

That premise, (C$_{\mathrm{loc}}$), has two parts: an entropy $\int\eta\,dA$ per cut with $\eta$ independent of the field strength and matter, and modes on every local horizon that are jointly KMS for the local boost before matter crosses. Deriving either part from the series' own structure is still open.

\end{document}